% Figura: rete SD-QKD e struttura interna di un nodo (adattata da GS QKD 015, Fig. 1)
\begin{figure}[t]
\centering
\begin{tikzpicture}[
  comp/.style={draw,rounded corners,align=center,font=\scriptsize,
               minimum height=6mm},
  small/.style={draw,circle,fill=green!12,minimum size=9mm,font=\scriptsize,
                align=center}]
  % controller
  \node[box,fill=blue!8,minimum width=4cm] (ctl) at (-0.6,4.3) {SDN controller (SD-QNC)};
  % nodo dettagliato
  \node[comp,fill=yellow!20,minimum width=1.4cm] (m1) at (-2.2,0) {modulo\\QKD 1};
  \node[comp,fill=yellow!20,minimum width=1.4cm] (m2) at (-0.6,0) {modulo\\QKD 2};
  \node[comp,fill=yellow!20,minimum width=1.4cm] (m3) at (1.0,0) {modulo\\QKD 3};
  \node[comp,fill=green!12,minimum width=4.6cm] (kms) at (-0.6,1.0) {key manager (KMS, key store)};
  \node[comp,fill=blue!6,minimum width=1.9cm] (agent) at (-1.95,2.0) {SDN agent};
  \node[comp,fill=gray!10,minimum width=2.4cm] (app) at (0.75,2.0) {applicazioni\\(QKD 004/014)};
  \begin{scope}[on background layer]
    \node[draw,dashed,rounded corners,fill=gray!3,inner sep=6pt,
          fit=(m1)(m3)(kms)(agent)(app),
          label={[font=\scriptsize\itshape]below:nodo SD-QKD (perimetro sicuro)}] (node) {};
  \end{scope}
  % nodi remoti
  \node[small] (r1) at (4.6,2.6) {nodo\\B};
  \node[small] (r2) at (4.4,-0.2) {nodo\\C};
  \node[small] (r3) at (-5.2,0.6) {nodo\\D};
  % collegamenti
  \draw[thick] (m3.east) -- node[lbl,above]{PHYS} (r2);
  \draw[thick] (m1.west) -- (r3);
  \draw[thick] (r1) -- (r2);
  \draw[thick,dotted,red!70!black] (kms.east) -- node[pos=0.55,font=\scriptsize,fill=white,inner sep=1pt]{VIRT via C} (r1);
  \draw[<->] (m1) -- (m1|-kms.south);
  \draw[<->] (m2) -- (m2|-kms.south);
  \draw[<->] (m3) -- (m3|-kms.south);
  \draw[<->] (agent.south) -- (agent.south|-kms.north);
  \draw[<->] (app.south) -- (app.south|-kms.north);
  \draw[<->,dashed,blue!60!black] (agent.north) -- node[pos=0.5,left,font=\scriptsize]{GS QKD 015} (ctl);
  \draw[<->,dashed,blue!60!black] (r1) -- (ctl);
  \draw[<->,dashed,blue!60!black] (r3) -- (ctl);
\end{tikzpicture}
\caption[Rete SD-QKD e struttura interna di un nodo]{Una rete SD-QKD: i nodi sono collegati da link fisici (linee
continue) e virtuali (punteggiata) e ciascuno è controllato dal controller
tramite il proprio SDN agent (tratteggio). Il nodo in primo piano mostra le
componenti interne: i moduli QKD, esposti come interfacce, il key manager,
l'interfaccia verso le applicazioni e l'agent. Adattata
da~\cite[Fig.~1]{qkd015}.}
\label{fig:sdqkd}
\end{figure}
